\chapter{Data Generation and Training Methodology}
\label{ch:training}

\section{Simulation Data Factory}

\subsection{Scenario Generation}

We generate diverse training scenarios by sampling:
\begin{itemize}
    \item \textbf{Routes}: Varying length (50--500 km), grade profiles (flat to \SI{2}{\percent}), curvature (straight to tight curves)
    \item \textbf{Consists}: Varying length (10--150 cars), mass distribution, car types
    \item \textbf{Control plans}: Different strategies (constant speed, coasting, aggressive acceleration/braking)
\end{itemize}

% TODO: Add table showing scenario distribution statistics
% TODO: Include figure showing example route/consist/control combinations

\subsection{Domain Randomization}

To improve generalization, we apply domain randomization:
\begin{itemize}
    \item Mass uncertainty: $\pm 5\%$ variation in car masses
    \item Resistance coefficients: $\pm 10\%$ variation in Davis equation parameters
    \item Adhesion: Varying friction coefficients (dry, wet, low adhesion)
    \item Initial conditions: Random initial speeds and positions
\end{itemize}

% TODO: Quantify impact of domain randomization on OOD performance

\section{Dataset Splits and Out-of-Distribution Tests}

The dataset is split into:
\begin{itemize}
    \item \textbf{Train} (70\%): Standard scenarios for model learning
    \item \textbf{Validation} (15\%): Held-out scenarios for hyperparameter tuning
    \item \textbf{Test} (15\%): Final evaluation on in-distribution data
\end{itemize}

Additionally, we define OOD test sets:
\begin{itemize}
    \item \textbf{Long consists}: 150+ cars (beyond training range)
    \item \textbf{Steep grades}: \SI{>2.5}{\percent} (extreme conditions)
    \item \textbf{Low adhesion}: Friction coefficient $< 0.15$ (safety-critical)
\end{itemize}

% TODO: Report performance degradation on OOD sets
% TODO: Add table comparing in-distribution vs OOD metrics

\section{Training Losses}

\subsection{Quantile Loss}

For quantile regression, we use the pinball loss:
\begin{equation}
    \mathcal{L}_{\text{quantile}}(y, \hat{q}_p) = \max(p(y - \hat{q}_p), (p-1)(y - \hat{q}_p))
    \label{eq:pinball_loss}
\end{equation}
where $\hat{q}_p$ is the predicted $p$-quantile.

\subsection{Calibration Loss}

To encourage calibration, we minimize the difference between predicted and empirical quantile coverage:
\begin{equation}
    \mathcal{L}_{\text{cal}} = \sum_p \left| \frac{1}{N} \sum_{i=1}^N \mathbf{1}(y_i \leq \hat{q}_{p,i}) - p \right|
    \label{eq:calibration_loss}
\end{equation}

\subsection{Physics Consistency Loss}

Physics penalties are implemented as soft constraints:
\begin{equation}
    \mathcal{L}_{\text{phys}} = \lambda_{\text{energy}} \mathcal{L}_{\text{energy}} + \lambda_{\text{force}} \mathcal{L}_{\text{force}} + \lambda_{\text{bounds}} \mathcal{L}_{\text{bounds}}
    \label{eq:physics_loss}
\end{equation}

% TODO: Detail each physics penalty term
% TODO: Report hyperparameter values ($\lambda$ coefficients)

\section{Implementation Details}

\subsection{Model Architecture}

\begin{itemize}
    \item Transformer: 6 encoder layers, 8 attention heads, hidden dimension 512
    \item Physics layer: Reduced-order model with 5 state variables per car
    \item Total parameters: $\sim 15$M
\end{itemize}

% TODO: Add architecture diagram with layer details
% TODO: Ablation on model size (4 vs 6 vs 8 layers)

\subsection{Training Procedure}

\begin{itemize}
    \item Optimizer: AdamW with learning rate $10^{-4}$, weight decay $10^{-5}$
    \item Batch size: 32 scenarios per batch
    \item Training duration: 50 epochs (early stopping on validation loss)
    \item Hardware: 4x NVIDIA A100 GPUs, training time $\sim 12$ hours
\end{itemize}

% TODO: Learning rate schedule details
% TODO: Report convergence curves

\subsection{Training Stability}

We employ several techniques to stabilize training:
\begin{itemize}
    \item Gradient clipping (max norm 1.0)
    \item Learning rate warmup (1000 steps)
    \item Mixed precision training (FP16)
    \item Physics loss annealing (start with small $\lambda$, increase gradually)
\end{itemize}

% TODO: Compare training curves with/without stability techniques
